\documentclass[10pt,a4paper]{amsart}
\usepackage[margin=19mm]{geometry}
\usepackage{amsmath,amssymb,mathtools}
\usepackage[scr=rsfs,cal=euler]{mathalfa}
\usepackage{xfrac}
\usepackage[unicode,hidelinks,bookmarks=false]{hyperref}
\newcommand{\ie}{i.e.\ }
\newcommand{\cf}{cf.\ }
\newcommand{\resp}{respectively\ }
\newcommand{\Subsection}{\subsection}
\providecommand{\nobreakdash}{\nobreak\hbox{-}}
\setlength{\emergencystretch}{2em}
\multlinegap0pt
\usepackage{bm}
\usepackage{bbm}
\usepackage{dsfont}
\usepackage{xspace}
\usepackage{cancel}
\usepackage{mathtools}
\usepackage{amsthm}
\newtheorem{theorem}{Theorem}
\newtheorem{lemma}[theorem]{Lemma}
\title[A Proof of Liu's Conjecture on the Fundamental Triangle Ineq]{A Proof of Liu's Conjecture on the Fundamental Triangle Inequality}
\author{Tserendorj Batbold}
\date{Source: arXiv:2609.10633v2 (CC BY 4.0)}
\begin{document}
\maketitle

\noindent\emph{Source and licence.} A Proof of Liu's Conjecture on the Fundamental Triangle Inequality. Version arXiv:2609.10633v2 (CC BY 4.0), distributed under the Creative Commons Attribution 4.0 International licence, as recorded in the arXiv licence metadata for this exact version and shown on the abstract page https://arxiv.org/abs/2609.10633v2. Full rights notice: licenses/arxiv-2609.10633-CC-BY-4.0.txt.\par\smallskip

\noindent\textit{Editorial scope.} Counted unit: the Lemma whose source label is \texttt{lem:power} in the article (source0.tex lines 101--117, source label \texttt{lem:power}), delivered together with the article's own complete original proof of it (source0.tex lines 119--236, one \texttt{proof} environment of 118 lines). No mathematics is written, repaired or re-derived: statement and proof are the article's own text line for line. Required context retained uncounted: none. Every definition, display and label of those lines is retained unchanged. Omitted from this package and not redistributed: 1--100, 118--118, 237--334. Nothing inside a retained excerpt is omitted or altered: every retained body is delivered exactly as the article's lines are written, so no omission receipt of any kind accompanies this package. Citations: the retained excerpts carry no citation command at all, so no bibliography entry of the article is transcribed and no reference is redistributed. Reference adaptation: none was needed and none was made; every reference inside the delivered unit resolves inside it, so no unresolved reference or citation is produced. Printed slips kept literal: no printed word, symbol, hypothesis or number of the counted statement or of its proof was altered. Layout and class: the article's own class is \texttt{amsart} with packages including geometry, amsmath, amssymb, amsthm, mathtools, microtype, enumitem, hyperref; the wrapper is the lane's pinned \texttt{amsart} editorial wrapper at 10pt on A4, which restates the theorem-like environments the retained text needs and adds the article's own macro definitions that the retained text needs (none), with extra wrapper packages (bm, bbm, dsfont, xspace, cancel, mathtools). The delivered numbering is the unit's own. Assets: no retained span contains a figure environment, a graphics-inclusion command, a PDF-page inclusion, a table or a TikZ picture; no image file is copied into this package, the package declares no asset dependency and no image file is redistributed with it. Licence: version arXiv:2609.10633v2 of this article is distributed under the Creative Commons Attribution 4.0 International licence, as recorded in the arXiv licence metadata for this exact version and shown on the abstract page https://arxiv.org/abs/2609.10633v2. A full rights notice is delivered at licenses/arxiv-2609.10633-CC-BY-4.0.txt. Characters: every retained excerpt is pure ASCII, so the delivered engine, XeLaTeX with its default Latin Modern text font, reports no missing character.
\begin{lemma}\label{lem:power}
Let $0<q\le 1$ and let $x,y,z>0$ satisfy
\begin{equation}\label{eq:constraints}
 x+y+z=2+q,\qquad xyz=q.
\end{equation}
Then
\begin{equation}\label{eq:power}
 x^k+y^k+z^k
 \begin{cases}
 \ge 2+q^k, & k>1,\\[2mm]
 \le 2+q^k, & k<1.
 \end{cases}
\end{equation}
If $k\ne0,1$, equality holds if and only if $(x,y,z)$ is a permutation of
$(q,1,1)$. For $k=0$ and $k=1$, equality holds for every
$(x,y,z)$ satisfying \eqref{eq:constraints}.
\end{lemma}

\begin{proof}
For $k=0$ and $k=1$, the assertion follows immediately from
\eqref{eq:constraints}. Hence assume $k\ne0,1$.

If $q=1$, then $x+y+z=3$ and $xyz=1$. Equality in AM--GM gives
$x=y=z=1$, and the assertion follows. Hence assume $0<q<1$.

Consider
\[
 \mathcal S_q=\{(x,y,z)\in(0,\infty)^3:x+y+z=2+q,\ xyz=q\}.
\]
This set is nonempty, since $(q,1,1)\in\mathcal S_q$, and it is compact. Indeed, every coordinate is at most $2+q$. Moreover,
\[
 yz\le \left(\frac{y+z}{2}\right)^2\le\frac{(2+q)^2}{4},
\]
so
\[
 x=\frac{q}{yz}\ge\frac{4q}{(2+q)^2},
\]
and cyclically the same positive lower bound holds for $y$ and $z$. Hence $\mathcal S_q$ is a closed subset of a compact box. Thus the continuous function
\[
 F(x,y,z)=x^k+y^k+z^k
\]
attains its minimum and maximum on $\mathcal S_q$.

We next verify that the constraint gradients are independent on $\mathcal S_q$. If
\[
 (1,1,1)\quad\text{and}\quad(yz,zx,xy)
\]
were linearly dependent, positivity would imply $x=y=z$. Then \eqref{eq:constraints} would yield
\[
 \left(\frac{2+q}{3}\right)^3=q,
\]
or
\[
 (2+q)^3-27q=(q-1)^2(q+8)=0,
\]
contrary to $0<q<1$. Hence the Lagrange multiplier conditions apply at every extremum of $F$ on $\mathcal S_q$.

At an extremum there exist real $\lambda,\mu$ such that
\[
 kx^{k-1}=\lambda+\mu yz,
\]
and cyclically. Since $xyz=q$, multiplication by the corresponding variable gives
\[
 kx^k-\lambda x-\mu q=0,
\]
with the corresponding equations for $y$ and $z$. Hence $x,y,z$ are positive zeros of
\[
 \phi(t)=kt^k-\lambda t-\mu q.
\]
For $t>0$,
\[
 \phi''(t)=k^2(k-1)t^{k-2},
\]
so $\phi'$ is strictly monotone. Hence $\phi$ has at most two positive
zeros: three distinct positive zeros would, by Rolle's theorem, give two
distinct zeros of $\phi'$, which is impossible. Therefore at least two
of $x,y,z$ are equal at every extremum.

By symmetry, write $y=z=t$. From \eqref{eq:constraints},
\[
 x=\frac{q}{t^2},\qquad \frac{q}{t^2}+2t=2+q.
\]
Hence
\[
 2t^3-(2+q)t^2+q=0,
\]
which factors as
\begin{equation}\label{eq:factor}
 (t-1)(2t^2-qt-q)=0.
\end{equation}

If $t=1$, then $(x,y,z)=(q,1,1)$ and
\begin{equation}\label{eq:F1}
 F_1=2+q^k.
\end{equation}

For the other positive solution of \eqref{eq:factor},
\[
 q=\frac{2t^2}{1+t},\qquad x=\frac{2}{1+t}.
\]
Since $0<q<1$, this positive solution satisfies $0<t<1$. The corresponding value of $F$ is
\[
 F_2=\left(\frac{2}{1+t}\right)^k+2t^k.
\]
Using $q=2t^2/(1+t)$, we obtain
\begin{align*}
 F_2-F_1
 &=\left(\frac{2}{1+t}\right)^k+2t^k-2
   -\left(\frac{2t^2}{1+t}\right)^k\\
 &=(1-t^k)\left[\left(\frac{2}{1+t}\right)^k(1+t^k)-2\right].
\end{align*}
For $k>1$, strict convexity of $u\mapsto u^k$ gives
\[
 \frac{1+t^k}{2}>\left(\frac{1+t}{2}\right)^k.
\]
Since $1-t^k>0$, it follows that $F_2>F_1$. Thus the global minimum is
$F_1$, proving the first inequality in \eqref{eq:power}.

If $0<k<1$, strict concavity gives
\[
 \frac{1+t^k}{2}<\left(\frac{1+t}{2}\right)^k.
\]
Since $1-t^k>0$, we obtain $F_2<F_1$.

If $k<0$, the function $u\mapsto u^k$ is strictly convex, and hence
\[
 \frac{1+t^k}{2}>\left(\frac{1+t}{2}\right)^k.
\]
In this case $1-t^k<0$, so again $F_2<F_1$. Therefore $F_1$ is the
global maximum for every $k<1$, $k\ne0$, and the second inequality in
\eqref{eq:power} follows.

The strictness of the above comparisons for $0<t<1$ shows that, for
$k\ne0,1$, equality can occur only when $(x,y,z)$ is a permutation of
$(q,1,1)$. Conversely, every permutation of $(q,1,1)$ gives equality.
\end{proof}

\end{document}
